% !TEX root = ../../main.tex
% C3.3 - Multi-object tracking (translated). Matches config/bytetrack_tracker.json and botsort_tracker.json.
\section{Multi-Object Tracking}
\label{sec:3.3}

\subsection{Problem and Basic Components}
\label{sec:3.3.1}
\label{sec:3.3.2}

A detector treats frames independently and cannot tell whether a person is the one seen before. Multi-Object Tracking (MOT) gives each object an identifier kept across frames; the sequence of its boxes is a \emph{track}, which in the project becomes an \emph{appearance}, the unit of search results (Section~\ref{sec:3.1.6}).

The common approach is \emph{tracking-by-detection}, founded on the SORT algorithm~\cite{bewley2016sort}, which has three components. First, each track has a Kalman filter~\cite{kalman1960} that predicts the position and size of its box in the current frame under a near-constant-velocity assumption, then corrects the prediction with the actual detection. Second, the match between a track and a detection is measured by the IoU (Equation~\ref{eq:iou}) between the predicted box and the detected box. Third, assigning each detection to at most one track at minimum total cost is solved with the Hungarian algorithm~\cite{kuhn1955hungarian}; an unmatched detection may start a new track, and an unmatched track is considered lost.

\subsection{ByteTrack}
\label{sec:3.3.3}

Earlier methods discarded low-confidence detections, so partially occluded people often broke their tracks. ByteTrack~\cite{zhang2022bytetrack} associates twice: high-confidence detections are matched with all tracks first, then the remaining tracks with low-confidence detections, so occluded people keep their tracks; unmatched low-confidence detections are dropped as mostly background. Only high-confidence detections start tracks, and a lost track is kept briefly so it can resume.

\subsection{BoT-SORT}
\label{sec:3.3.4}

BoT-SORT~\cite{aharon2022botsort} extends ByteTrack by adjusting the Kalman filter state, compensating camera motion and adding appearance re-identification (ReID) features to the association step. In the project, BoT-SORT is an alternative tracker with the same thresholds as ByteTrack; camera motion compensation and ReID are disabled because the cameras are static and to keep the cost low on a CPU-only machine.

\subsection{Track Life Cycle in the System}
\label{sec:3.3.5}

The tracker receives the detections of the sampled frames, separately for each camera, with the following rules:
\begin{itemize}
    \item \textbf{Confidence thresholds:} detections of 0.25 or more are high and may start a new track; detections from 0.1 to below 0.25 are only used in the second association step.
    \item \textbf{Confirmation:} a track is confirmed only after at least two observations; an unconfirmed track produces no result.
    \item \textbf{Termination:} a track ends after more than three lost processed frames or two seconds without observation, or when the stream ends. Both video files and RTSP streams use the timestamps recorded in the stream itself (Section~\ref{sec:3.7.4}), so the rule measures the time that passed in front of the camera and does not depend on how fast the machine processes.
    \item \textbf{Identifier scope:} one camera and one processing session only; the system does not link the tracks of the same person across cameras.
\end{itemize}
